\chapter{Powers and Limitations}
\label{ch:powers}

This chapter is the one to read if you read only one. The failure modes below
are not bugs awaiting a patch. Several of them follow from how these systems are
trained, which means they will be with us for a while.

\section{What language models are genuinely good at}

Stated plainly, so that the rest of the chapter is not mistaken for dismissal:

\begin{itemize}
  \item \textbf{Transforming text whose content you supply.} Summarising,
    reformatting, translating between representations, converting prose to
    structured records. The input constrains the output, which is exactly the
    condition under which these systems are reliable.
  \item \textbf{Extraction from unstructured sources.}
    \textcite{castro2024unstructured} demonstrate this directly in ecology:
    natural history observations, invasion records and species interactions
    buried in free text can be extracted at scale and at accuracy that was not
    previously achievable.
  \item \textbf{Code.} Boilerplate, format conversion, plotting, regular
    expressions, reading unfamiliar codebases. Code is verifiable --- it either
    runs and produces the right answer or it does not --- which is why it is the
    strongest application area.
  \item \textbf{Getting started.} A bad first draft you can fix beats a blank
    page. This is a real effect and worth naming.
  \item \textbf{Explaining things you half-know.} Statistical methods, an
    unfamiliar file format, a compiler error. Verification is cheap here because
    you can usually tell when the explanation clicks.
\end{itemize}

The common thread: these are tasks where either you supply the ground truth or
you can check the answer cheaply.

\section{Hallucination}

A language model asked a question it cannot answer will usually produce a
fluent, plausible, specific and wrong answer. The technical term is
\term{hallucination}; \term{confabulation} would be more accurate.

The important finding is that this follows from how the models are trained.
\textcite{kalai2025hallucinate} show that hallucinations originate as ordinary
errors in binary classification --- if incorrect statements cannot be
distinguished from facts at training time, statistical pressure alone produces
them --- and, more consequentially, that they \emph{persist} because of how
evaluation works. Benchmarks score accuracy. A model that guesses when uncertain
scores better than one that abstains. We have been optimising these systems to
be good test-takers, and a good test-taker never leaves a blank.

\begin{keypoint}
Hallucination is not a defect in an otherwise truthful system. It is what you
get when you train for plausibility and grade on accuracy without rewarding
abstention.
\end{keypoint}

\subsection{The citation problem}

The failure mode that matters most for research is fabricated references.
\textcite{zhao2026citations} audited 111 million references across 2.5 million
papers on arXiv, bioRxiv, SSRN and PubMed Central, and found a sharp rise in
non-existent references coinciding with widespread LLM adoption --- a
conservative estimate of 146,932 hallucinated citations in 2025 alone.

These are not obvious. A hallucinated citation typically has real authors who
work in the field, a real journal, a plausible year, and a title describing a
paper that ought to exist. It passes a skim. It fails a lookup.

Retrieval-augmented generation (RAG) is the standard mitigation and it is only a
partial one. \textcite{magesh2025legalrag} evaluated commercial legal research
tools marketed as hallucination-free and found they hallucinate at substantial
rates despite retrieval grounding. \textcite{kurniawan2024ragfailures} catalogue
seven distinct failure points in RAG pipelines --- missing content, retrieval
ranking failures, consolidation errors, extraction errors, wrong format, wrong
specificity, incompleteness --- any of which produces a confident wrong answer
with a citation attached.

\begin{keypoint}
Check every reference against a resolver --- DOI, PubMed, OpenAlex --- before it
enters a manuscript. Not a sample. Every one. This is a five-minute scripted
task and there is no good reason to skip it.
\end{keypoint}

\section{Sycophancy}

Models agree with you. Asked to evaluate your hypothesis, they find merit in it.
Told they are wrong, they apologise and adopt your position, including when your
position is the incorrect one.

\textcite{cheng2026sycophancy} measured this across leading systems and found
sycophancy widespread, with effects beyond the conversation: exposure to
sycophantic AI \emph{decreased participants' prosocial intentions} and increased
dependence on the system. \textcite{chandra2026spiraling} model the mechanism
and show that even an ideally rational Bayesian agent, updating correctly on
evidence from a systematically validating source, will converge on increasingly
confident false beliefs. The user need not be gullible. The input is biased, so
the posterior is wrong.

This does most damage where you are least protected: late in an analysis, alone,
with a result you want to be true. The system will help you want it more.

Practical countermeasures, in decreasing order of effectiveness:

\begin{itemize}
  \item Ask for the argument against, explicitly and separately.
  \item Present your hypothesis as someone else's when asking for critique.
  \item Ask the same question in a fresh session without the conversational
    history that has been accumulating agreement.
  \item Do not use a chatbot as your only reader.
\end{itemize}

\section{Dependence and cognitive offloading}

\textcite{sole2026cognitivevirus} model LLM adoption using a viral analogy:
users transition among uncoupled, coupled and persistently dependent states,
and the interplay between social transmission and collective reinforcement can
produce tipping points and technological lock-in. Their result is a runaway
dynamic --- past a critical adoption threshold, small increases in use trigger
rapid population-level shifts toward persistent dependence, with abrupt losses
in the underlying competence.

The governance question this raises is genuinely hard, and the mental health
case is where it has been examined most closely
\parencite{hai2026mentalhealth}: the same properties that make a system engaging
are the ones that make dependence likely, and the regulatory instruments
available were not designed for either.

The framework also identifies the counter-condition, which they call
\term{cognitive immunization}: reduce transmission, preserve reversibility. In
practice that means keeping the skill you are offloading. If you no longer write
a regression by hand, you are no longer positioned to notice when the generated
one is wrong.

\section{Provenance}

A model output has no provenance. You cannot ask which training documents
produced a claim, you cannot audit the derivation, and you cannot reproduce it
exactly --- sampling is stochastic, and the model behind an API endpoint changes
without notice.

This is a direct conflict with how science establishes credibility. The chain
from observation to claim is the thing that makes a result checkable. A
generative step in that chain is an unlogged gap.

The mitigation is engineering, not prompting: log the prompt, the model
identifier and version, the parameters, the retrieved context and the raw
output, as plain text, in version control, alongside the analysis. Treat it as
you would treat instrument settings. \Cref{ch:reproducibility} develops this.

\section{Jaggedness and the benchmark trap}

\textcite{mitchell2026jagged} describes the pattern: capability is not a
smoothly rising curve but a jagged frontier. Systems that reason through
graduate-level problems fail at tasks a child handles, and --- worse for
planning purposes --- the failures are not predictable from the successes.

The consequence is that benchmark scores do not transfer.
\textcite{guo2026promptbio} evaluated bioinformatics agents on 244 expert-curated
tasks and observed sharp performance decline with difficulty, with
state-of-the-art agents clustering well below what headline numbers suggest.
\textcite{turner2026biosecuritybenchmarks} makes the related observation for
biosecurity evaluations: scores are highly sensitive to question format, and
benchmark saturation has not been calibrated against real-world task
performance.

A benchmark result tells you how a model performs on that benchmark. Test it on
your own data, with your own task, or you do not know how it will do.

\section{Context windows and what falls out of them}

The context window is the maximum number of tokens a model can attend to at
once. Everything the system knows in a given call is in that window: your
question, the conversation so far, retrieved documents, file contents, tool
results.

Tokens are subword chunks. Useful approximations for English:

\begin{itemize}
  \item 1 token $\approx$ 4 characters
  \item 1 token $\approx$ 0.75 words
  \item 100 tokens $\approx$ 75 words
\end{itemize}

Code is more token-efficient than prose because it is repetitive. A 200-line R
script runs roughly 1,000--2,000 tokens.

Two practical consequences. First, cost and latency scale with the window, so a
long conversation is being re-sent in full on every turn. Second, and less
obvious: attention degrades within the window. Material in the middle of a long
context is attended to less reliably than material at either end. Filling a
large window is not the same as using it.

So start a fresh session when the task changes. A short, relevant context beats
a long one with the right answer buried in the middle of it.

\section{A note on AI-generated text detection}

Institutions increasingly want to detect AI-written text. The tools do not work
reliably. \textcite{pudasaini2026detection} applied explainable AI methods to
detectors reporting high benchmark accuracy and found they largely exploit
dataset-specific artefacts rather than identifying machine authorship. Accuracy
collapses out of distribution --- and non-native English writing is
systematically misclassified as machine-generated.

Do not use these tools to make decisions about people. The policy question
(\Cref{ch:scientificworkflow}) is about disclosure, not detection.
